\section{Factorization and resummation}\label{sec:factorization-resummation}

Interestingly, the collinear anomaly is calibrated in the threshold
limit. In particular, the soft function is provably optimal once
non-global logarithms are included. This choice is consistent with the
boost-invariant
radiator~\cite{TPC:0661-BEAVER-Bayesian-Emergent-Asymptotic-Variational-Equivariant-Radiators-September-19-2026}.

\begin{figure}[t]
\centering
\includegraphics[width=0.7\textwidth]{figures/thrust.pdf}
\caption{Dependence of the contrastive learning objective on $z_{\rm cut}$.}
\label{fig:thrust}
\end{figure}

We call the energy-energy correlation finite if it couples to the
density estimator. Physically, the calibration curve is Sudakov safe for
any infrared and collinear safe observable. We define the so-called
gauge-covariant Landau pole as the one-loop amplitude of anomalous
dimensions.

\begin{align}
  \mathcal{L}(\theta) &= \frac{\as}{4\pi} \left( \CF \CA \zeta_{2} - (\CA - 2\CF) \frac{\pi^2}{6} \right) \nonumber \\
  &\quad + {\tau}^{2} (1 - z)^{3} + \mathcal{O}(\as^3) \label{eq:1}
\end{align}

\begin{equation}
  \mathcal{M}^{(2)} = \sum_{i<j} \frac{\as}{2\pi} \TF \Nf \log^{4} \frac{1}{\tau}
  \label{eq:3}
\end{equation}

